\documentclass[11pt]{article}
\usepackage{amsmath,amssymb,amsthm,booktabs}
\usepackage[margin=1in]{geometry}
\newtheorem{theorem}{Theorem}
\newtheorem{lemma}[theorem]{Lemma}
\newtheorem{definition}[theorem]{Definition}
\title{Problem 877: XOR-Equation A}
\author{Project Euler Solutions}
\date{}
\begin{document}
\maketitle

\section{Problem Statement}

The XOR-product $x \otimes y$ is carry-less multiplication in base~2. Consider:
\[
(a \otimes a) \oplus (2 \otimes a \otimes b) \oplus (b \otimes b) = 5.
\]
Let $X(N)$ be the XOR of all $b$ for solutions with $0 \le a \le b \le N$. Given $X(10)=5$, find $X(10^{18})$.

\section{Mathematical Foundation}

\begin{definition}[$\operatorname{GF}(2)$ Polynomial Ring]
$n = \sum_i b_i 2^i \longleftrightarrow n(x) = \sum_i b_i x^i \in \operatorname{GF}(2)[x]$. XOR $=$ addition; XOR-product $=$ multiplication; left-shift by 1 $=$ multiplication by $x$.
\end{definition}

\begin{theorem}[Frobenius Endomorphism]
In $\operatorname{GF}(2)[x]$: $a(x)^2 = a(x^2)$.
\end{theorem}

\begin{proof}
In characteristic~2, $(c+d)^2 = c^2+d^2$. By induction, $(\sum_i b_i x^i)^2 = \sum_i b_i x^{2i} = a(x^2)$ since $b_i^2 = b_i$ for $b_i \in \{0,1\}$.
\end{proof}

\begin{theorem}[Equation Reformulation]
The equation becomes, in $\operatorname{GF}(2)[x]$:
\[
a(x^2) + x \cdot a(x) \cdot b(x) + b(x^2) = k(x),
\]
where $k=5$ gives $k(x) = x^2 + 1$.
\end{theorem}

\begin{proof}
$a \otimes a \leftrightarrow a(x)^2 = a(x^2)$ by Frobenius. $2 \otimes a \otimes b \leftrightarrow x \cdot a(x) \cdot b(x)$ since $2 \leftrightarrow x$. Similarly for $b \otimes b$. Addition in $\operatorname{GF}(2)[x]$ is XOR.
\end{proof}

\begin{lemma}[Bit Decomposition]
Writing $a(x) = a_0(x^2) + x\,a_1(x^2)$, $b(x) = b_0(x^2) + x\,b_1(x^2)$, the equation separates into constraints on even- and odd-degree coefficients of $k(x)$.
\end{lemma}

\begin{proof}
Substitute and expand. The $x^4$ terms and cross-products separate by parity, yielding a system in $a_0, a_1, b_0, b_1$ solvable recursively.
\end{proof}

\begin{theorem}[Recursive Structure]
Solutions have a binary-tree structure: lowest bits are determined by lowest-degree coefficients of $k(x)$, and remaining bits satisfy a reduced equation. This yields $O(\log N)$-depth enumeration.
\end{theorem}

\begin{proof}
After fixing parity bits to match constant and linear terms, higher-order terms satisfy an analogous equation with halved degree. The process terminates in $O(\log N)$ steps.
\end{proof}

\section{Algorithm}

Digit DP on binary representations of $a$ and $b$. At each bit position, the $\operatorname{GF}(2)$ constraints determine valid bit choices. Track tightness (whether $b \le N$) and ordering ($a \le b$).

\section{Complexity Analysis}

\begin{itemize}
\item \textbf{Time:} $O(\log^2 N)$. The DP has $O(\log N)$ levels with $O(1)$ branching per state.
\item \textbf{Space:} $O(\log N)$ for recursion stack.
\end{itemize}

\section{Answer}

\[
\boxed{336785000760344621}
\]

\end{document}
